\section{结论与展望}\label{sec:conclusion}

DD1、DD2 分别在线性与一般函数尺度下，给出了次单调性和残差正则性的定量补偿关系。
DD3 用有限维全局条件保证近端步存在，参数分析则将补偿条件用于选择固定或变化的近端参数。
在各自的存在性、留域及长度条件下，可得到相应收敛结论。

一般函数下，距离下降、点收敛与有限长度之间仍缺少完整刻画。
本文的比较级数和积分条件是充分条件；实际长度还可能取决于相邻步的方向。
进一步的应用需要从具体算子的结构中推导缺陷与正则性函数。
